\section{Coefficient extraction from the generating functions}
\label{app:coefficient-extraction}

Let
\begin{equation}
  G(x,y,t)=\sum_{i,j\geq0}p_{i,j}(t)x^iy^j
  \label{eq:appendix-pgf-series}
\end{equation}
be analytic in a neighbourhood of the origin.  The bivariate Taylor formula gives
\begin{equation}
  p_{i,j}(t)
  =[x^iy^j]G(x,y,t)
  =\frac{1}{i!\,j!}
   \left.
   \frac{\partial^{i+j}G}
        {\partial x^i\partial y^j}
   \right|_{x=y=0}.
  \label{eq:appendix-state-probability}
\end{equation}
For completeness, the coefficient is also represented by the bivariate Cauchy
integral
\begin{equation}
  p_{i,j}(t)
  =\frac{1}{(2\pi\mathrm i)^2}
   \int_{|x|=r_x}\int_{|y|=r_y}
   \frac{G(x,y,t)}
        {x^{i+1}y^{j+1}}\,\dd y\,\dd x,
  \label{eq:bivariate-cauchy}
\end{equation}
for circles inside the domain of analyticity.  This form is useful for stable
numerical extraction when derivatives of a special-function representation are
unwieldy.

\subsection{A multilinear power}

Suppose
\begin{equation}
  G(x,y,t)=\{c(t)+f(t)x+g(t)y\}^{N}.
  \label{eq:multilinear-power}
\end{equation}
For \(a+b\leq N\), repeated differentiation gives
\begin{equation}
  \frac{\partial^{a+b}G}{\partial x^a\partial y^b}
  =
  \frac{N!}{(N-a-b)!}
  f^ag^b(c+fx+gy)^{N-a-b},
  \label{eq:multilinear-derivative}
\end{equation}
and the derivative is zero when \(a+b>N\).  Hence
\begin{equation}
  p_{i,j}(t)
  =
  \frac{N!}{i!\,j!\,(N-i-j)!}
  f(t)^i g(t)^j c(t)^{N-i-j},
  \qquad i+j\leq N.
  \label{eq:trinomial-probability}
\end{equation}
This is the trinomial law.  In particular, the factor brought down by
differentiation is the falling factorial \(N!/(N-i-j)!\), not \(N^{i+j}\).
The distinction is essential for exact normalisation.

For transfer alone,
\begin{equation}
  c=0,\qquad
  f=1-\e^{-\alpha t},\qquad
  g=\e^{-\alpha t}.
\end{equation}
Only states with \(i+j=N\) have positive probability, and
\begin{equation}
  p_{i,j}(t)
  =\binom Ni
   (1-\e^{-\alpha t})^i
   \e^{-\alpha tj},
  \qquad i+j=N.
  \label{eq:transfer-only-coefficients}
\end{equation}

For transfer followed by death, take
\begin{align}
  c(t)
  &=1-\frac{\mu}{\mu-\alpha}\e^{-\alpha t}
       +\frac{\alpha}{\mu-\alpha}\e^{-\mu t},
  \notag\\
  f(t)
  &=\frac{\alpha}{\mu-\alpha}
    (\e^{-\alpha t}-\e^{-\mu t}),
  \notag\\
  g(t)&=\e^{-\alpha t},
  \label{eq:transfer-death-cfg}
\end{align}
with the continuous equal-rate limit from
\cref{eq:transfer-death-equal-rates}.  Substitution into
\cref{eq:trinomial-probability} gives the complete joint law.  Since
\(c+f+g=1\), summing over \(i\) and \(j\) recovers one by the multinomial theorem.

\subsection{Transfer, birth and death}

The full generating function is not a finite multilinear power in \(x\), but its
dependence on \(y\) remains linear before the power is taken.  From
\cref{eq:abd-h-definition},
\begin{equation}
  G(x,y,t)
  =\{\e^{-\alpha t}y+h(x,t)\}^{N}.
\end{equation}
The binomial theorem gives
\begin{equation}
  G(x,y,t)
  =\sum_{j=0}^{N}
   \binom Nj
   \e^{-\alpha tj}y^j
   h(x,t)^{N-j},
\end{equation}
and therefore
\begin{equation}
  p_{i,j}(t)
  =\binom Nj\e^{-\alpha tj}
   [x^i]h(x,t)^{N-j}.
  \label{eq:appendix-abd-coefficients}
\end{equation}

If
\begin{equation}
  h(x,t)=\sum_{k\geq0}h_k(t)x^k,
\end{equation}
then the remaining coefficient is the finite convolution
\begin{equation}
  [x^i]h(x,t)^M
  =
  \sum_{\substack{k_1+\cdots+k_M=i\\k_r\geq0}}
  \prod_{r=1}^{M}h_{k_r}(t).
  \label{eq:coefficient-convolution}
\end{equation}
Truncating \(h\) after degree \(i\) is exact for this calculation: coefficients of
higher degree cannot contribute to \([x^i]\).  This observation underlies
\texttt{scripts/extract\_coefficients.py}.  As a check, summing over all interior
states at fixed \(j\) means setting \(x=1\); since
\(h(1,t)=1-\e^{-\alpha t}\),
\begin{equation}
  \sum_{i=0}^{\infty}p_{i,j}(t)
  =
  \binom Nj
  \e^{-\alpha tj}
  (1-\e^{-\alpha t})^{N-j},
\end{equation}
which is the exterior binomial marginal in
\cref{eq:abd-exterior-binomial}.

